% ============================================================
% Supplemental Material for
%   "Exact nonequilibrium steady state of collective group-flip dynamics
%    on quenched interaction neighborhoods"
% REVTeX 4.2, one column. Overleaf: compile this file separately
% (Menu -> Main document -> SM.tex). Needs refs_pre.bib and
% fig4_controlled_test.pdf (or .png).
% ============================================================
\documentclass[aps,pre,onecolumn,superscriptaddress,amsmath,amssymb,longbibliography]{revtex4-2}

\usepackage{graphicx}
\usepackage{amsthm}
\usepackage{bm}
\usepackage[hyphens]{url}
\usepackage[colorlinks=true,linkcolor=blue,citecolor=blue,urlcolor=blue]{hyperref}

\renewcommand{\thesection}{S\arabic{section}}
\renewcommand{\theequation}{S\arabic{equation}}
\renewcommand{\thetable}{S\arabic{table}}
\renewcommand{\thefigure}{S\arabic{figure}}

\newtheorem{theorem}{Theorem}[section]
\newtheorem{lemma}[theorem]{Lemma}
\newtheorem{proposition}[theorem]{Proposition}
\newtheorem{corollary}[theorem]{Corollary}
\newtheorem{claim}[theorem]{Claim}
\newtheorem{conjecture}[theorem]{Conjecture}
\newtheorem{openq}[theorem]{Open problem}
\newtheorem{heuristic}[theorem]{Heuristic}
\theoremstyle{definition}
\newtheorem{protocol}[theorem]{Protocol}
\newtheorem{note}[theorem]{Note}
\theoremstyle{remark}
\newtheorem{remark}[theorem]{Remark}

\newcommand{\ctil}{\tilde c}
\newcommand{\cG}{c_G}
\newcommand{\E}{\mathbb{E}}

\begin{document}

\title{Supplemental Material for ``Exact nonequilibrium steady state of
collective group-flip dynamics on quenched interaction neighborhoods''}

\author{Pham Le Dat Nguyen}
\email{phmldtngyn.777@gmail.com}
\affiliation{Independent Researcher, Ha Noi, Viet Nam}

\maketitle

% ============================================================
\section{Notation, conventions and correspondence}
\label{sec:S1}
% ============================================================

\paragraph*{Objects.} $G=(V,E)$ is a connected graph with $|V|=n$ and
graph distance $d$, which is symmetric. Ball $B_i=B_i(r)=\{j:d(i,j)\le r\}$,
$\beta_i=|B_i|$, and $\beta$ when all balls have the same size; coverage
$q=\beta/n$; residual $R_i=V\setminus B_i$. Conditions: \textbf{R1}
$\beta_i>n/2$ for all $i$; \textbf{R2} $\mathrm{diam}(R_i)\le r$ for all
$i$. Dynamics: each node fires at rate one; a fire at $i$ moves all nodes
of state $\sigma_i$ inside $B_i$ to a target state drawn uniformly among
the other $Q-1$ states ($Q=2$: the other state). Observables: stationary
majority fraction $M$, magnetization $m$ with $|m|=2M-1$; number $K$ of
updates between consecutive rotations (changes of the majority label)
and cycle time $T$; containment observables $\cG$, $\ctil$ as in the main
text.

\paragraph*{Graphs used.} The ring $C_n(1,2)$ (links to first and second
neighbors on a cycle) and the two-dimensional lattice with first- and
second-neighbor links along each axis are the 2-section graphs of the
uniform hypergrids whose hyperedges are three consecutive nodes along an
axis. The group-flip dynamics on a hypergraph depends on it only through
its 2-section graph, since $B_i(r)$ is the graph-metric ball of radius $r$
there. All statements below are therefore stated for graphs.

\paragraph*{Status of statements.} \emph{Theorem, Proposition, Lemma,
Corollary}: complete proof given. \emph{Claim}: proof sketch only; not
used as a hypothesis of any result. \emph{Conjecture}: numerical evidence
only. \emph{Heuristic}: structural statement without proof.
\emph{Protocol}: measurement or statistical design. \emph{Note}:
methodological record.

\paragraph*{Correspondence with the main text.} Main-text
Theorem~1 is Proposition~\ref{prop:rot}; Corollary~1 is
Corollary~\ref{cor:dual}; Proposition~1 is Proposition~\ref{prop:irrev};
Lemmas~1--2 are Lemmas~\ref{lem:epoch}--\ref{lem:overlap}; Theorem~2 is
Theorem~\ref{thm:band}; Corollary~2 combines Theorem~\ref{thm:cycle} and
Corollary~\ref{cor:obs}; Proposition~2 is Proposition~\ref{prop:het};
Proposition~3 is Propositions~\ref{prop:E1}--\ref{prop:E2};
Conjecture~1 is Conjecture~\ref{conj:upper}.

% ============================================================
\section{Well-mixed limit and multistate version}
\label{sec:S2}
% ============================================================

\begin{lemma}[Count-chain reduction]
\label{lem:A1}
If $B_i=V$ for all $i$, the state-count vector $x=(x_1,\dots,x_Q)$,
$\sum_kx_k=n$, is a Markov chain with transitions
$x\to x+x_k(e_l-e_k)$ at rate $x_k/(Q-1)$ for each ordered pair $(k,l)$,
$k\ne l$.
\end{lemma}
\begin{proof}
A fire at a node of state $k$ (total rate $x_k$) moves the entire class
$k$ to the target $l$, chosen with probability $1/(Q-1)$. The rates
depend on the configuration only through $x$.
\end{proof}

\begin{lemma}[Pure death of live states]
\label{lem:A2}
The number $L_t$ of states with at least one node is a pure-death chain
on $\{1,\dots,Q\}$ with $L\to L-1$ at rate $n(L-1)/(Q-1)$; $L=1$ is
absorbing for $L$.
\end{lemma}
\begin{proof}
A fire from live class $k$ decreases $L$ if and only if the target is
another live class, which has probability $(L-1)/(Q-1)$; summing over
live $k$ gives rate $n(L-1)/(Q-1)$. A target that is an absent state only
relabels the firing class.
\end{proof}

\begin{theorem}[First-consensus time]
\label{thm:A3}
If $L_0$ states are initially live,
\[
\E[T_{\rm cons}]=\frac{Q-1}{n}\sum_{\ell=2}^{L_0}\frac{1}{\ell-1}=\frac{Q-1}{n}H_{L_0-1};
\]
in particular $\E[T_{\rm cons}]=1/n$ for $Q=2$.
\end{theorem}
\begin{proof}
Sum the mean holding times of Lemma~\ref{lem:A2}.
\end{proof}

\begin{theorem}[Minority advantage, $Q=2$]
\label{thm:A4}
From $(a,n-a)$, the class of size $a$ wins the first consensus with
probability $1-a/n$.
\end{theorem}
\begin{proof}
The first fire decides: if it is in the class of size $a$ (probability
$a/n$), that class moves into the other state and loses.
\end{proof}

\begin{theorem}[Three states, one absent]
\label{thm:A5}
Let $Q=3$, states 1 and 2 live with sizes $a,b>0$, $a+b=n$, and state 3
absent. With $D=a^2+3ab+b^2$, the probabilities that state $k$ wins the
first consensus are $P_1=b(a+b)/D$, $P_2=a(a+b)/D$, $P_3=ab/D$.
\end{theorem}
\begin{proof}
Let $F$ be the probability that state 1 wins when states 1 and 2 are
live with sizes $a$ and $b$, $G$ the same with the roles of $a$ and $b$
exchanged, and $U$ the probability that state 1 wins when it is absent
and the live classes have sizes $a$ and $b$. First-step analysis gives
\begin{align*}
(a+b)F&=a\big[\tfrac12\cdot0+\tfrac12U\big]+b\big[\tfrac12\cdot1+\tfrac12F\big],\\
(a+b)U&=a\big[\tfrac12F+\tfrac12\cdot0\big]+b\big[\tfrac12G+\tfrac12\cdot0\big],\\
(a+b)G&=b\big[\tfrac12\cdot0+\tfrac12U\big]+a\big[\tfrac12\cdot1+\tfrac12G\big].
\end{align*}
The first and third equations give $F=(aU+b)/(2a+b)$ and
$G=(bU+a)/(2b+a)$. Substituting into the second, with $t=a+b$, $p=ab$ and
$(2a+b)(2b+a)=2t^2+p$, gives $U=p/(t^2+p)=ab/D$; back-substitution gives
$F$ and $G$.
\end{proof}

\begin{theorem}[General $Q$]
\label{thm:A6}
Consensus states are absorbing in the well-mixed limit and every other
state is transient; the winning probabilities are the unique solution of
$(I-P_{TT})x=P_{TA}x_A$, where $P_{TT}$ is the substochastic transient
block, invertible by transience. No closed form is claimed for $Q\ge4$.
\end{theorem}

% ============================================================
\section{Ball geometry}
\label{sec:S3}
% ============================================================

\begin{proposition}[Ball sizes without wrap-around]
\label{prop:B1}
On the $d$-dimensional lattice with first- and second-neighbor links
along each axis, the cost of an offset $o$ along an axis is
$\lceil|o|/2\rceil$ and, without wrap-around,
$|B_r|=\sum_{c_1+\dots+c_d\le r}\prod_kw(c_k)$ with $w(0)=1$, $w(c)=4$ for
$c\ge1$. Hence $|B_r|=4r+1$ in one dimension and $8r^2+1$ in two.
\end{proposition}
\begin{proof}
Offsets of cost $c\ge1$ along an axis are $\pm(2c-1),\pm2c$; convolve over
axes. In 2D, $(1+4r)+\sum_{c=1}^{r}4(1+4(r-c))=8r^2+1$.
\end{proof}

\begin{proposition}[Ring: distance, ball size, R1, R2]
\label{prop:B2}
On $C_n(1,2)$, $d(i,j)=\lceil\delta(i,j)/2\rceil$ and $\beta=\min(n,4r+1)$.
For every finite $n$: R1 $\iff2\beta>n$, and R2 $\iff\beta=n$ or
$3\beta\ge2n-1$; when $\beta=4r+1<n$ the latter is $n\le6r+2$.
\end{proposition}
\begin{proof}
A step moves $\pm1$ or $\pm2$ sites, so balls are arcs. For $\beta<n$ the
residual is an arc of $n-\beta$ nodes, of graph diameter
$\lceil(n-\beta-1)/2\rceil$, computed along the arc because
$n-\beta<n/2$ under R1. R2 reads $n-\beta-1\le2r$; with $r=(\beta-1)/4$
this is $3\beta\ge2n-1$.
\end{proof}

\begin{proposition}[Coverage radius]
\label{prop:B3}
For fixed $0<q<1$ and $d$, the smallest radius with $|B_r|\ge qn$ in the
no-wrap regime satisfies $r_{\min}=\Theta(n^{1/d})$.
\end{proposition}
\begin{proof}
$|B_r|$ is a strictly increasing polynomial of degree $d$ in $r$.
\end{proof}

\begin{claim}[Wrapped regime; sketch only]
\label{cl:B4}
For fixed $q$ on $d$-dimensional tori, $|B_r|/n$ converges to the
rescaled volume of the corresponding diamond with $O(n^{-1/d})$
corrections. Checked numerically only; not used below.
\end{claim}

\begin{claim}[Sparse configuration model; sketch only]
\label{cl:B5}
If $(k-1)^{2r+1}=o(n)$, then $|B_r|$ in the configuration model is
asymptotically the non-backtracking tree volume. Not used below; it
motivates the truncated-ball design of Protocol~\ref{prot:reg}.
\end{claim}

% ============================================================
\section{Proofs of the exact results}
\label{sec:S4}
% ============================================================

\subsection{Results valid under R1}

\begin{proposition}[Rotation identity; equal balls not required]
\label{prop:rot}
Let $Q=2$ and $\beta_i>n/2$ for every $i$. From a strict-majority state
with majority $\mathcal M$ and minority $\mathcal A$, a rotation occurs
if and only if the firing node lies in $\mathcal M$; the rotation rate is
$|\mathcal M|$. Tie states are transient and carry zero stationary mass.
\end{proposition}
\begin{proof}
Majority fire at $i$: the class sizes after the fire are
$|\mathcal A|+|\mathcal M\cap B_i|$ (new state) and
$|\mathcal M\setminus B_i|$ (old state), with difference
$(2\beta_i-n)+2|\mathcal A\setminus B_i|>0$. Minority fire at $i$:
$\mathcal M\cup(\mathcal A\cap B_i)$ grows and $\mathcal A\setminus B_i$
shrinks. Ties: the same computation with $|\mathcal M|=|\mathcal A|$
shows that any fire produces a strict majority, and strict-majority
states map to strict-majority states.
\end{proof}

\begin{remark}
\label{rem:onlymaj}
In a majority fire only majority nodes inside $B_i$ move; minority nodes
inside $B_i$ already carry the target state. Only $\beta_i>n/2$ for the
firing node is used.
\end{remark}

\begin{corollary}[Static--dynamic identity]
\label{cor:dual}
Under R1, on each closed class $c$ of the configuration chain with
stationary law $\mu_c$, the long-run fraction of fires that are rotations
equals $\E_{\mu_c}[M]$, the long-run rotation rate per unit time is
$n\E_{\mu_c}[M]$, and the long-run mean number of fires per rotation is
$1/\E_{\mu_c}[M]$.
\end{corollary}
\begin{proof}
Write the dynamics as a Poisson clock of rate $n$ and, at each ring, a
firing node uniform in $V$ and independent of the past. The embedded
jump chain $P=n^{-1}\sum_iT_i$, with $T_i$ the deterministic update of a
fire at $i$, has the same stationary law as the continuous-time chain
because the total rate is constant. By the ergodic theorem for the jump
chain restricted to $c$, the fraction of fires at which the configuration
$x$ has firing node in its majority converges to
$\sum_x\mu_c(x)|\mathcal M(x)|/n=\E_{\mu_c}[M]$, and by
Proposition~\ref{prop:rot} these are exactly the rotations. Fires occur
at rate $n$, which gives the rotation rate, and the mean number of fires
per rotation is the reciprocal of the rotation fraction.
\end{proof}

\begin{proposition}[One-way erasure transitions]
\label{prop:irrev}
Assume R1. Let $x$ have nonempty minority $\mathcal A$, and let a fire of
$j\in\mathcal A$ with $\mathcal A\subseteq B_j$ map $x$ to the consensus
configuration $y$. Then there is no $k$ with $T_k(y)=x$. If $\mu$ is a
stationary law with $\mu(x)>0$, then $\mu$ is not reversible.
\end{proposition}
\begin{proof}
From consensus, $T_k$ switches exactly $B_k$; so $T_k(y)$ has its
nonconsensus state on $B_k$, whereas $x$ has it on $\mathcal A$, and
$|\mathcal A|<n/2<\beta_k$. Reversibility,
$\mu(x)w(x\to y)=\mu(y)w(y\to x)$, would require $w(y\to x)>0$ because
$\mu(x)w(x\to y)\ge\mu(x)\cdot1>0$.
\end{proof}

\begin{remark}
\label{rem:currents}
The stationary rate of minority fires is $n(1-\E[M])$ by the argument of
Corollary~\ref{cor:dual}. Under R1 and R2 every minority fire is an
erasure (Lemma~\ref{lem:epoch}), so the total stationary probability
current through one-way transitions is $n(1-\E[M])>0$ whenever
$\E[M]<1$. In Region~II only the fraction of minority fires that erase
the minority is one-way in the sense above; the irreversibility
conclusion holds as soon as one such transition has positive stationary
rate; this is the case for the six odd Region-II cases of
Table~\ref{tab:smexact}, whose reachable set is a single irreducible
class containing the consensus states, with $\ctil>0$.
\end{remark}

\subsection{Exact steady state under R1 and R2}

\begin{lemma}[Entry and initial conditions]
\label{lem:entry}
Let $Q=2$ and assume R1$\wedge$R2. (i)~From consensus, a fire at $i$
creates the minority $R_i$. (ii)~From any strict-majority state, the
minority after the next rotation is contained in $R_i$, where $i$ is the
node whose fire produced it. (iii)~Ties are left at the first update and
never re-entered. (iv)~From any initial configuration the epoch chain
enters a closed class in finite time almost surely. Different closed
classes may carry different stationary laws; Lemma~\ref{lem:xbar} shows
that the stationary mean minority fraction is the same on all of them.
\end{lemma}
\begin{proof}
(i)~The moved class is $B_i$, which is the new majority by R1. (ii)~is
Lemma~\ref{lem:epoch}. (iii)~is Proposition~\ref{prop:rot}. (iv)~The
state space is finite.
\end{proof}

\begin{lemma}[Epoch invariant and self-containment]
\label{lem:epoch}
Under R1$\wedge$R2 with $Q=2$, every post-rotation minority $A_m$,
created by a fire at $i_m$, satisfies $A_m\subseteq R_{i_m}$. Hence
$A_m\subseteq B_j$ for all $j\in A_m$ and $|A_m|\le n-\beta_{i_m}$.
Equal balls are not required.
\end{lemma}
\begin{proof}
Base case: Lemma~\ref{lem:entry}(i). Let $A\subseteq R_i$ be the current
minority. R2 gives $R_i\subseteq B_j$ for every $j\in R_i$, hence
$A\subseteq B_j$ for every $j\in A$. A fire at $j\in A$ moves the whole
minority, the system reaches consensus, and the following fire at some
$i'$ creates $A'=R_{i'}$. A fire at $i\in\mathcal M$ is a rotation and the
new minority is $\mathcal M\setminus B_i\subseteq R_i$.
\end{proof}

\begin{lemma}[Overlap identity; equal balls]
\label{lem:overlap}
Under Lemma~\ref{lem:epoch} with $\beta_i\equiv\beta$, for $i$ uniform in
$\mathcal M$: $\E[|A\cap B_i|\mid a]=a(\beta-a)/(n-a)$.
\end{lemma}
\begin{proof}
$\sum_{i\in\mathcal M}|A\cap B_i|=\sum_{j\in A}|B_j\cap\mathcal M|$ by
symmetry of $d$. For $j\in A$, $B_j\supseteq A$, so
$|B_j\cap\mathcal M|=\beta-a$.
\end{proof}

\begin{lemma}[Affine conditional mean]
\label{lem:affine}
Under the assumptions of Lemma~\ref{lem:overlap},
$\E[a_{m+1}\mid a_m=a]=(1-q)(n-a)$.
\end{lemma}
\begin{proof}
Minority fire (probability $a/n$): $a'=n-\beta$. Majority fire
(probability $1-a/n$): $a'=n-\beta-a+|A\cap B_i|$, with conditional mean
$(n-\beta)(n-2a)/(n-a)$. Weighting gives $(n-\beta)(n-a)/n$.
\end{proof}

\begin{lemma}[Stationary epoch mean and cycle structure]
\label{lem:xbar}
For any invariant distribution $\pi$ of the epoch chain,
$\bar x=\E_\pi[a]/n=(1-q)/(2-q)$. Moreover $P(K=2\mid a)=a/n$ and
$K\in\{1,2\}$, so $\E_\pi[K]=1+\bar x$.
\end{lemma}
\begin{proof}
Take $\pi$-expectations in Lemma~\ref{lem:affine}. A cycle contains a
second fire if and only if its first fire is a minority fire; after the
resulting consensus, the next fire is a rotation.
\end{proof}

\begin{remark}
\label{rem:emptyminority}
A post-rotation minority can be empty: for instance, right after a fire
at $i$ from consensus ($A=R_i$), a second fire at $i$ restores
consensus. Then $a=0$, $P(K=2\mid a)=0$ and the first interval has
$M=1$, consistently with the formulas above.
\end{remark}

\begin{theorem}[Exact stationary majority fraction]
\label{thm:band}
Assume $Q=2$, R1, R2 and equal balls $\beta$. Then
$\E[M]=(2-q)/(3-2q)$, for every initial law.
\end{theorem}
\begin{proof}
Start the epoch chain in an invariant $\pi$. Let $Y_m=\int_{\text{cycle }m}M\,dt$
and $\tau_m$ be the reward and length of cycle $m$; $(Y_m,\tau_m)$ is
stationary, not independent. By the Birkhoff ergodic theorem,
$N^{-1}\sum_{m\le N}Y_m\to\E[Y\mid\mathcal I]$ and
$N^{-1}\sum_{m\le N}\tau_m\to\E[\tau\mid\mathcal I]\ge1/n$ almost
surely. On the event that the chain is in closed class $c$, these are
the class means. Inter-update times are exponential with mean $1/n$ and
independent of the sequence of firing nodes. During the first interval
$M=1-a_m/n$; if the first fire is a minority fire (probability $a_m/n$),
the second interval has $M=1$. Hence $\E_c[Y]=n^{-1}(1-\bar x_c)+n^{-1}\bar x_c=1/n$
and $\E_c[\tau]=(1+\bar x_c)/n$, with $\bar x_c=(1-q)/(2-q)$ by
Lemma~\ref{lem:xbar} applied to the chain restricted to $c$. The ratio
is $1/(1+\bar x)=(2-q)/(3-2q)$ on every class, hence almost surely; by
Lemma~\ref{lem:entry}(iv) the same holds from any initial law.
\end{proof}

\begin{theorem}[Two-point cycle law]
\label{thm:cycle}
Under the assumptions of Theorem~\ref{thm:band}: $\E[K]=1+\bar x=1/\E[M]$,
$\E[T]=(1+\bar x)/n$, and
$\mathrm{CV}_T^2=(7-8q+2q^2)/(3-2q)^2$.
\end{theorem}
\begin{proof}
$T=G_1+BG_2$ with $G_{1,2}$ exponential of mean $1/n$ and $B$ Bernoulli
with mean $\bar x$; $B$ is determined by the firing node at the end of
the first interval, hence independent of $G_1,G_2$. Then
$\mathrm{Var}(T)=n^{-2}(1+2\bar x-\bar x^2)$; with
$\bar x=(1-q)/(2-q)$, $1+\bar x=(3-2q)/(2-q)$ and
$1+2\bar x-\bar x^2=(7-8q+2q^2)/(2-q)^2$.
\end{proof}

\begin{corollary}[Further closed forms]
\label{cor:obs}
Under the assumptions of Theorem~\ref{thm:band}:
$\E|m|=2\E[M]-1=1/(3-2q)$; long-run rotation rate
$n(2-q)/(3-2q)$ (Corollary~\ref{cor:dual}); long-run rate of one-way
erasure transitions $n(1-\E[M])=n(1-q)/(3-2q)$
(Remark~\ref{rem:currents}).
\end{corollary}

\begin{corollary}[Threshold location of the closed form]
\label{cor:thr}
$f(q)=(2-q)/(3-2q)$ is strictly increasing, $f'(q)=(3-2q)^{-2}$,
$f(1/2)=3/4$, $f(2/3)=4/5$. A crossing of $M=3/4$ predicted by $f$ would
lie at $q=1/2$, outside the exact regime; this concerns the closed form,
not the dynamics there. For two-dimensional lattices R1$\wedge$R2 was
verified by breadth-first search at the sizes used ($L=12$, $r=4$,
$\beta=111$; $L=30$, $r=10$, $\beta=699$).
\end{corollary}

\subsection{Unequal balls}

\begin{proposition}[Unequal balls]
\label{prop:het}
Assume $Q=2$, R1, R2, and $\beta_i<n$ for some $i$. Let
$\bar\beta=n^{-1}\sum_i\beta_i$, $\bar q=\bar\beta/n$,
$S_A=\sum_{j\in A}\beta_j$, $q_{\min}=\min_i\beta_i/n$,
$q_{\max}=\max_i\beta_i/n$. For each closed class $c$ of the epoch chain
with invariant law $\pi_c$, let
$q^\ast_c=\E_{\pi_c}[S_A]/(n\E_{\pi_c}[a])$. Then
$q^\ast_c\in[q_{\min},q_{\max}]$,
\begin{equation}
\E_c[M]=g(q^\ast_c),\qquad g(s)=\frac{2+\bar q-2s}{3-2s},
\label{eq:Shet}
\end{equation}
and $g(q_{\max})\le\E[M]\le g(q_{\min})$ for every initial law.
\end{proposition}
\begin{proof}
\emph{Conditional mean.} Lemma~\ref{lem:epoch} holds without equal
balls. Given the minority $A$ of size $a$: a minority fire (probability
$a/n$) erases $A$, and the next fire, at a uniform node $i'$, creates
$R_{i'}$ of mean size $n-\bar\beta$. A majority fire at $i$ uniform in
$\mathcal M$ creates the minority $\mathcal M\setminus B_i$, of size
$(n-a)-\beta_i+|A\cap B_i|$. Now
$\sum_{i\in\mathcal M}\beta_i=n\bar\beta-S_A$ and, by symmetry of $d$ and
self-containment,
$\sum_{i\in\mathcal M}|A\cap B_i|=\sum_{j\in A}|B_j\cap\mathcal M|=\sum_{j\in A}(\beta_j-a)=S_A-a^2$.
Hence
\[
\E[a'\mid A,\text{maj.}]=n-a+\frac{2S_A-n\bar\beta-a^2}{n-a},
\]
and weighting the branches,
\[
\E[a'\mid A]=\frac an(n-\bar\beta)+\frac{n-a}{n}\Big(n-a+\frac{2S_A-n\bar\beta-a^2}{n-a}\Big)
=n-\bar\beta-a\Big(1+\frac{\bar\beta}{n}\Big)+\frac{2S_A}{n}.
\]
\emph{Stationary mean.} Under $\pi_c$,
$\E[a](2+\bar q)=n(1-\bar q)+2\E[S_A]/n=n(1-\bar q)+2q^\ast_c\E[a]$, so
$\bar x_c=(1-\bar q)/(2+\bar q-2q^\ast_c)$. Since $\beta_i<n$ for some
$i$, a fire at $i$ from consensus creates a nonempty minority, so
$\E_{\pi_c}[a]>0$ and $q^\ast_c$ is well defined;
$a\,n q_{\min}\le S_A\le a\,nq_{\max}$ gives the range.
\emph{Renewal--reward.} The cycle structure is unchanged
($K\in\{1,2\}$, $P(K=2\mid a)=a/n$), so the proof of
Theorem~\ref{thm:band} gives $\E_c[M]=1/(1+\bar x_c)=g(q^\ast_c)$.
\emph{Bounds.} $g'(s)=-2(1-\bar q)/(3-2s)^2\le0$, so
$g(q_{\max})\le g(q^\ast_c)\le g(q_{\min})$ for every class, hence for the
almost-sure long-run average from any initial law.
\end{proof}

\begin{remark}
\label{rem:het-numerics}
\emph{Numerical verification} (\path{scripts/hetero_balls_check.py}).
We generated rings with $9\le n\le14$ plus up to five random chords,
radius $r\in\{2,3,4\}$, kept the ball systems with unequal sizes, no full
ball, and R1$\wedge$R2 verified by breadth-first search, and removed
duplicate ball systems; 60 distinct systems remained. For each, the
stationary law of the configuration chain on the states reachable from
consensus was obtained by a sparse direct solve ($L^1$ residual
$\le3.1\times10^{-15}$). The epoch law was computed independently as the
rotation flux into each state, $\propto\sum_x\mu(x)n^{-1}\#\{i\in\mathcal
M(x):T_i(x)=y\}$, and $q^\ast$ from it. Results: $\max|\E[M]-g(q^\ast)|=1.2\times10^{-15}$;
the bounds hold in all 60 cases; $\max|(\text{rotation fraction})-\E[M]|=2.0\times10^{-14}$
(Corollary~\ref{cor:dual}); $q^\ast<\bar q$ in all 60 cases, and
$\E[M]-(2-\bar q)/(3-2\bar q)\in[0.0008,0.0127]$. No proof of
$q^\ast\le\bar q$ is known.
\end{remark}

% ============================================================
\section{Containment theory}
\label{sec:S5}
% ============================================================

\begin{proposition}[$\cG$ on vertex-transitive and configuration geometries]
\label{prop:D1}
$\cG=\E_{i,j}\mathbf 1\{R_i\subseteq B_j\}$ with $i,j$ independent and
uniform. On a vertex-transitive graph with uniform $\beta$,
$\cG=n^{-1}|\{j:R_i\subseteq B_j\}|$, independent of $i$. In general
$\cG$ is not determined by ball sizes and must be computed by
enumeration.
\end{proposition}

\begin{proposition}[Necessary condition]
\label{prop:D2}
If $R_i\subseteq B_j$ then $\mathrm{diam}(R_i)\le2r$.
\end{proposition}
\begin{proof}
$d(u,v)\le d(u,j)+d(j,v)\le2r$ for $u,v\in R_i\subseteq B_j$.
\end{proof}

\begin{proposition}[R2 implies one-step extinction]
\label{prop:D3}
Under R2, $\{j\in R_i:R_i\subseteq B_j\}=R_i$ for every $i$, hence
$\ctil=1$. Uniform balls are not required.
\end{proposition}

\begin{proposition}[Intersection form, equal balls]
\label{prop:D4}
$R_i\subseteq B_j\iff B_i\cup B_j=V\iff|B_i\cap B_j|=2\beta-n$; in
particular $\cG=0$ if $2\beta<n$. Hence $\cG$ is computed exactly by
counting pairs with $\mathrm{popcount}(B_i\,\&\,B_j)=2\beta-n$.
\end{proposition}

\begin{proposition}[$\cG$ on the ring]
\label{prop:E1}
With $\beta=\min(n,4r+1)$, $\cG=\min\{1,\max(0,2\beta-n+1)/n\}$; at
$2\beta=n$, $\cG=1/n$.
\end{proposition}
\begin{proof}
For $\beta<n$, $\beta$ is odd and a ball arc of length $\beta$ contains
the residual arc of length $n-\beta$ if and only if its center lies in an
interval of $2\beta-n+1$ positions; divide by $n$. For $\beta=n$
residuals are empty.
\end{proof}

\begin{proposition}[$\ctil$ on the ring]
\label{prop:E2}
$\ctil=0$ if $2\beta<n$; $\ctil=(2\beta-n+1)/(n-\beta)$ if $2\beta\ge n$
and $3\beta<2n-1$ (equal to $1/\beta$ at $2\beta=n$); $\ctil=1$ if
$3\beta\ge2n-1$. It jumps from $0$ to $1/\beta$ at $2\beta=n$ and is
continuous at $3\beta=2n-1$.
\end{proposition}
\begin{proof}
Intersect the interval of admissible centers with the residual. At
$2\beta=n$ containment forces $B_j=R$, whose unique center lies in $R$.
If $2\beta>n$ and $3\beta<2n-1$ the interval has length
$2\beta-n+1<n-\beta$ and lies inside $R$; if $3\beta\ge2n-1$ it covers
$R$.
\end{proof}

\begin{proposition}[Three regions on the ring]
\label{prop:E3}
Region~I: $3\beta\ge2n-1$ (R1 and R2, $\ctil=1$). Region~II:
$n/2<\beta<(2n-1)/3$ (R1 only, $0<\ctil<1$). Boundary: $2\beta=n$
($\cG=1/n$, $\ctil=1/\beta$). Region~III: $2\beta<n$ ($\ctil=0$).
\end{proposition}

\begin{proposition}[Newborn identity]
\label{prop:newborn}
Let $Q=2$. From consensus, a fire at $i$ creates the minority $R_i$.
Conditional on the next event being a fire of a minority node, the
firing node is uniform in $R_i$, and the minority is erased if and only
if $R_i\subseteq B_j$. Averaging over $i$ with weights $|R_i|$ gives
$\ctil$.
\end{proposition}

\paragraph*{Validation.} Propositions~\ref{prop:E1}--\ref{prop:E2} agree
exactly (rational arithmetic, explicit set inclusion) with exhaustive
enumeration for every $n\in[12,40]$ and every radius, 192 cases,
including the boundary case $(n,r)=(18,2)$ with $\cG=1/18$, $\ctil=1/9$
(\path{tests/independent/containment_bruteforce.py}). The Monte Carlo
estimate of the newborn erasure probability agrees with $\ctil$ within
binomial error on native ring balls, e.g., $n=15$, $r=2$: 29\,489 trials,
$0.6664\pm0.0027$ against $2/3$; $n=25$, $r=3$: 19\,535 trials,
$0.1666\pm0.0027$ against $1/6$; and in all 18 cells with $0<\ctil<1$ of
Protocol~\ref{prot:reg} (largest $|z|=1.85$).

% ============================================================
\section{Exact solvers, simulations and annealed computations}
\label{sec:S6}
% ============================================================

\begin{protocol}[Exact solvers and simulation checks]
\label{prot:exact}
\emph{(i) Reduced exact solver} (\path{tests/gvm_validation_suite.py},
task T5). States are orbits of reachable minority sets under the
dihedral group and state exchange; lazy power iteration from the
consensus orbit, $L^1$ residual $\le1.4\times10^{-14}$
(Table~\ref{tab:smexact}).
\emph{(ii) Independent cross-check}
(\path{tests/independent/exact_chain_independent.py}). Full
configurations, no symmetry reduction, transitions generated literally
from the update rule, irreducibility of the reachable set checked; sparse
direct solve ($n\le21$) or lazy power iteration to $L^1$ change
$<10^{-16}$ ($n=23,25$). The lumped chain is also solved in exact
rational arithmetic (FLINT, Dixon's algorithm) and verified by exact
substitution (Table~\ref{tab:smcheck}).
\emph{(iii) Simulations} (\path{tests/independent/sim_independent.py}).
Independent continuous-time simulator, 10 replicas of $2\times10^6$
fires per case. Region~II, native ring balls, the six cases of
Table~\ref{tab:smexact}: largest $|M_{\rm sim}-M_{\rm exact}|=4.6\times10^{-5}$,
largest $|z|=1.6$. Region~I, native balls with R1 and R2 verified by
breadth-first search (ring $n=400$, $r=75$; $n=401$, $r=70$; 2D $L=12$,
$r=4$; $L=30$, $r=10$): largest $|M_{\rm sim}-f(q)|=2.1\times10^{-5}$
(largest $|z|=0.8$); cycle law (Theorem~\ref{thm:cycle}): largest
$|z|=1.9$ for $\E[K]$ and for $\mathrm{CV}_T^2$.
\end{protocol}

\begin{protocol}[Extended Region-II solutions and annealed version]
\label{prot:ext}
\emph{(i) Quenched chain} (\path{scripts/region2_exact.c}). Written
independently of Protocol~\ref{prot:exact}. States: configurations
reachable from consensus by breadth-first search, identified under state
exchange only. Since every node fires at rate one, the stationary law
equals that of the jump chain $P=n^{-1}\sum_iT_i$; power iteration of
$(I+P)/2$ until the $L^1$ change is below $10^{-12}$, then 3000 further
iterations; residual $\|\mu P-\mu\|_1<5\times10^{-13}$ in every case. It
reproduces the six exact rationals to $10^{-12}$.
\emph{(ii) Annealed version} (\path{scripts/annealed_baseline.py}). At
every fire of $i$ the interaction set is $\{i\}$ plus $\beta-1$ nodes
drawn uniformly without replacement from $V\setminus\{i\}$. If $i$ has a
state shared by $m$ nodes, $1+h$ nodes switch, with $h$ hypergeometric
(population $n-1$, $m-1$ successes, $\beta-1$ draws). The chain on the
count $n_1$ is solved by a direct linear solve (residual
$<2\times10^{-13}$ up to $n=400$, $<10^{-12}$ up to $n=3200$).
\emph{(iii) Independent re-implementation}
(\path{scripts/annealed_size_scan.py}). A separate implementation of (ii)
reproduces all 17 annealed values of main-text Table~II to
$5\times10^{-7}$ (the table's rounding) and the four $n=400$ values
quoted in the main text, and produces Table~\ref{tab:sizescan} and the
$n=400$ curve of main-text Fig.~2 (\path{scripts/fig_quenched_annealed.py}).
\end{protocol}

\begin{table}[t]
\caption{Annealed version: exact stationary $\E[M]_{\rm ann}$ for
$\beta=\mathrm{round}(qn)$ and the gap $f(\beta/n)-\E[M]_{\rm ann}$ to the
quenched exact value. Last row: large-$n$ random map of
Heuristic~\ref{heur:map} ($4\times10^6$ iterations, batch-means standard
error).}
\label{tab:sizescan}
\begin{ruledtabular}
\begin{tabular}{rcccccc}
 & \multicolumn{2}{c}{$q=0.7$ ($f=0.81250$)} & \multicolumn{2}{c}{$q=0.8$ ($f=0.85714$)} & \multicolumn{2}{c}{$q=0.9$ ($f=0.91667$)}\\
$n$ & $\E[M]_{\rm ann}$ & gap & $\E[M]_{\rm ann}$ & gap & $\E[M]_{\rm ann}$ & gap\\
\colrule
100  & 0.79524 & 0.01726 & 0.85031 & 0.00683 & 0.91550 & 0.00116\\
200  & 0.79499 & 0.01751 & 0.85016 & 0.00699 & 0.91544 & 0.00122\\
400  & 0.79486 & 0.01764 & 0.85008 & 0.00706 & 0.91541 & 0.00125\\
800  & 0.79480 & 0.01770 & 0.85004 & 0.00710 & 0.91540 & 0.00127\\
1600 & 0.79477 & 0.01773 & 0.85002 & 0.00712 & 0.91539 & 0.00127\\
3200 & 0.79475 & 0.01775 & ---     & ---     & ---     & ---\\
map  & $0.79472(2)$ & --- & $0.84999(2)$ & --- & $0.91539(1)$ & ---\\
\end{tabular}
\end{ruledtabular}
\end{table}

At $n=400$ the annealed value is below $f(\beta/n)$ for every
$\beta\le n-2$ in Region~I of the ring ($267\le\beta\le398$), with gaps
from $1.2\times10^{-7}$ ($\beta=398$) to $0.022$; at $\beta=n-1$ the gap
is zero to machine precision, since each interaction set then misses a
single node.

\begin{heuristic}[Large-$n$ annealed map]
\label{heur:map}
For $q>1/2$ and $n\to\infty$ at fixed $q$, hypergeometric concentration
suggests that the minority fraction $x$ of the annealed version evolves,
per fire, by the random map $x\to(1-q)x$ with probability $x$ (minority
fire) and $x\to(1-q)(1-x)$ with probability $1-x$ (majority fire, always
a rotation for $q>1/2$). No proof of convergence is attempted; the
stationary mean of $1-x$ under this map agrees with the exact chain at
the largest sizes computed (Table~\ref{tab:sizescan}). The quenched
counterparts are exact: $x\to0$ for a minority fire
(Lemma~\ref{lem:epoch}) and conditional mean $(1-q)(1-2x)/(1-x)$ for a
majority fire (Lemma~\ref{lem:overlap}).
\end{heuristic}

\begin{table}[t]
\caption{Exact chain in Region~II of the ring, odd $n\le25$. Reachable:
minority sets reachable from consensus (one irreducible class in every
case); orbits: classes under the dihedral group and state exchange.}
\label{tab:smexact}
\begin{ruledtabular}
\begin{tabular}{rrrcccccrr}
$n$ & $r$ & $\beta$ & $q$ & $\cG$ & $\ctil$ & $\E[M]$ & $f(q)$ & reachable & orbits\\
\colrule
 9 & 1 &  5 & 0.5556 & 2/9  & 1/2 & 43/57   & 13/17 &    73 &    8\\
15 & 2 &  9 & 0.6000 & 4/15 & 2/3 & 65/84   & 7/9   &   481 &   24\\
17 & 2 &  9 & 0.5294 & 2/17 & 1/4 & 149/201 & 25/33 &  2177 &   80\\
21 & 3 & 13 & 0.6190 & 2/7  & 3/4 & 265/339 & 29/37 &  2689 &   80\\
23 & 3 & 13 & 0.5652 & 4/23 & 2/5 & 72/95   & 33/43 & 11777 &  288\\
25 & 3 & 13 & 0.5200 & 2/25 & 1/6 & 319/433 & 37/49 & 51201 & 1088\\
\end{tabular}
\end{ruledtabular}
\end{table}

\begin{table}[t]
\caption{Independent cross-check of Table~\ref{tab:smexact}: full
configurations, no symmetry reduction. Config.: reachable configurations.
Last column: floating-point value minus the exact rational.}
\label{tab:smcheck}
\begin{ruledtabular}
\begin{tabular}{rrrccc}
$n$ & $r$ & config. & $\E[M]$ (floating point) & $L^1$ residual & difference\\
\colrule
9 & 1 & 146 & 0.7543859649 & $1.4\times10^{-15}$ & $-1.1\times10^{-16}$\\
15 & 2 & 962 & 0.7738095238 & $1.3\times10^{-14}$ & $1.1\times10^{-16}$\\
17 & 2 & 4354 & 0.7412935323 & $2.7\times10^{-14}$ & $1.2\times10^{-14}$\\
21 & 3 & 5378 & 0.7817109145 & $4.1\times10^{-14}$ & $5.7\times10^{-15}$\\
23 & 3 & 23554 & 0.7578947368 & $9.6\times10^{-16}$ & $-2.1\times10^{-15}$\\
25 & 3 & 102402 & 0.7367205543 & $9.4\times10^{-15}$ & $2.1\times10^{-15}$\\
\end{tabular}
\end{ruledtabular}
\end{table}

% ============================================================
\section{Exploratory regression test of the one-step containment hypothesis}
\label{sec:S7}
% ============================================================

\begin{protocol}[Exploratory regression test]
\label{prot:reg}
\emph{Design.} Families: ring (cycle); two-dimensional torus with king
(eight-neighbor) adjacency; configuration model with degree sequence
$k=6$, drawn as a multigraph and projected to a simple graph (self-loops
removed, multi-edges merged), hence not exactly 6-regular;
Watts--Strogatz with $k=6$, $p=0.1$. Sizes $n\in\{100,144,400\}$;
coverages $q\in\{0.55,0.60,\dots,0.80\}$. One graph per (family, $n$),
seed 42; three runs (seeds 1, 2, 3) per (family, $n$, $q$): 216 runs.

\emph{Interaction sets.} $\beta=\mathrm{round}(qn)$; $B_i$ consists of the
$\beta$ nodes nearest to $i$ under breadth-first layering, ties in the
last layer broken at random with the run seed. The three runs of a cell
therefore share the graph but not necessarily the interaction sets;
$\ctil$ differs between runs in 15 of the 72 cells. Truncated sets need
not satisfy $j\in B_i\iff i\in B_j$, so this is a modified
interaction-set model, not the native-ball model of the theorems.

\emph{Measurement.} $Q=2$; $2\times10^5$ fires per run, $5\times10^4$
discarded; $M$ and its Monte Carlo standard error from 75 batch means
of 2000 fires; $\cG$ and $\ctil$ by exact enumeration on the truncated
sets. Data: \path{data/processed/measurements_v2.csv} (one row per run).

\emph{Model.} Weighted least squares, weights $1/\mathrm{se}^2$:
$M\sim q+C(\text{family})+C(n)$ (reduced) and
$M\sim q+\ctil+C(\text{family})+C(n)$ (full); $\gamma$ is the coefficient
of $\ctil$.

\emph{Family-effect norm.}
$N=[\sum_{f<g}(\alpha_f-\alpha_g)^2]^{1/2}$ over all six pairs of family
effects, invariant under the choice of reference; reduction
$(N_{\rm red}-N_{\rm full})/N_{\rm red}$.

\emph{Primary bootstrap.} Within each family, its three (family, $n$)
clusters are resampled with replacement; within each selected cluster,
the three runs are resampled inside each of the six coverage strata.
$B=2000$, seed 9999, percentile intervals.

\emph{Acceptance criteria (fixed before analysis).} (a)~$\gamma>0$ with
interval excluding zero; (b)~reduction $\ge0.30$.

\emph{Result.} $\hat\gamma=-0.0196$ [95\% interval $(-0.0222,-0.0169)$];
reduction $-0.98$ [$(-1.13,-0.85)$]; fraction of replicates with
reduction $\ge0.30$: $0.000$. Neither criterion is met.

\emph{Status.} Exploratory: the resampling scheme and the family
statistic were chosen after a first analysis (Note~\ref{note:record}).

\emph{Limitations.} (1)~$\ctil$ describes only the first minority fire
after a minority is created from consensus. (2)~Within families, $\ctil$
varies almost only in the ring family (96.6\% of the within-family sum
of squares; it is identically zero in all 2D cells and at most 0.241
(WS) and 0.036 (configuration)), so $\gamma$ is weakly identified.
(3)~Only 12 clusters, each a single graph realization. (4)~Coverage
enters linearly in the primary model. (5)~A failed criterion does not
identify an alternative mechanism.
\end{protocol}

\begin{note}[Analysis record and sensitivity analyses]
\label{note:record}
Only the acceptance criteria were fixed before the data were analyzed.
The initial analysis resampled runs independently and measured family
separation relative to the ring reference; an intermediate analysis used
a two-level cluster bootstrap and the norm of the three contrasts among
the non-ring families. Both were replaced, for the reasons below; all
items are reproduced by \path{src/analysis/reanalysis_v2.py}.
\emph{1. Point estimate.} Full-sample estimates: $-0.0191$ without
$C(n)$, $-0.0196$ with $C(n)$; the intermediate analysis reported the
bootstrap mean ($-0.0165$). All point estimates quoted are full-sample
fits.
\emph{2. Family statistic.} The three-contrast norm omits all contrasts
involving the ring family, the only one in which $\ctil$ varies
appreciably, and is not reference-invariant. With all six contrasts the
reduction is $-0.98$; the three-contrast statistic gives $-0.047$
[$(-0.167,+0.018)$].
\emph{3. Bootstrap.} The unstratified two-level scheme omitted at least
one family in 61 of 500 replicates (258 of 2000 in the primary
specification) and did not preserve coverage strata. All its replicates
with $\hat\gamma>0$ lack at least one family; its wide interval
[$(-0.0236,+0.0022)$ in the primary specification] reflects replicates
in which $\gamma$ is barely identified.
\emph{4. Standard errors.} The between-run standard deviation within a
cell divided by the mean batch-means standard error has median 0.832 over
the 72 cells, against $\sqrt{\ln2}\approx0.833$ expected for three runs
with correct standard errors. Fits to the 72 cell means give
$\hat\gamma=-0.0188$.
\emph{5. Dependence.} Resampling runs independently gives
$(-0.0223,-0.0170)$; resampling only within cells (Monte Carlo noise
alone) gives $(-0.0200,-0.0192)$.
\emph{6. Coverage variable.} Using the realized $\beta/n$ instead of the
nominal $q$ (difference $\le0.5/n$) gives $\hat\gamma=-0.0201$.
\end{note}

\begin{table}[t]
\caption{Regression test, all specifications and resampling schemes.
Point: full-sample estimate. Intervals: bootstrap percentile 95\%.
Reduction: six-contrast family-norm reduction (negative means increased
separation). Missing: replicates lacking at least one family. $\dagger$:
specification of the intermediate analysis (no $C(n)$, $B=500$).}
\label{tab:smreg}
\footnotesize\setlength{\tabcolsep}{3pt}
\begin{ruledtabular}
\begin{tabular}{llrcccr}
Specification & Bootstrap & $B$ & $\hat\gamma$ & $\gamma$ 95\% & reduction 95\% & missing\\
\colrule
lin.\ $q$, no $C(n)$$^\dagger$ & unstratified cluster & 500 & $-0.0191$ & $(-0.0230,+0.0045)$ & $(-1.269,-0.003)$ & 61\\
lin.\ $q$, no $C(n)$$^\dagger$ & runs iid & 500 & $-0.0191$ & $(-0.0220,-0.0163)$ & $(-1.251,-0.737)$ & 0\\
primary: lin.\ $q$, $C(n)$ & unstratified cluster & 2000 & $-0.0196$ & $(-0.0236,+0.0022)$ & $(-1.365,-0.001)$ & 258\\
primary: lin.\ $q$, $C(n)$ & stratified cluster & 2000 & $-0.0196$ & $(-0.0222,-0.0169)$ & $(-1.129,-0.849)$ & 0\\
primary: lin.\ $q$, $C(n)$ & runs iid & 2000 & $-0.0196$ & $(-0.0223,-0.0170)$ & $(-1.249,-0.753)$ & 0\\
primary: lin.\ $q$, $C(n)$ & within cells & 2000 & $-0.0196$ & $(-0.0200,-0.0192)$ & $(-1.021,-0.937)$ & 0\\
lin.\ $q$, no $C(n)$ & stratified cluster & 2000 & $-0.0191$ & $(-0.0221,-0.0165)$ & $(-1.116,-0.821)$ & 0\\
categ.\ $q$, $C(n)$ & stratified cluster & 2000 & $-0.0185$ & $(-0.0208,-0.0161)$ & $(-1.074,-0.802)$ & 0\\
$(n,q)$-cell FE & stratified cluster & 2000 & $-0.0186$ & $(-0.0215,-0.0159)$ & $(-1.119,-0.772)$ & 0\\
primary, OLS & --- & --- & $-0.0181$ & --- & point $-0.84$ & ---\\
primary, 72 cell means & --- & --- & $-0.0188$ & --- & point $-0.94$ & ---\\
\end{tabular}
\end{ruledtabular}
\end{table}

\begin{figure}[t]
\centering
\includegraphics[width=\textwidth]{fig4_controlled_test}
\caption{Exploratory test. (a)~Simulated $M$ against $q$ at $n=400$ for
the four families (mean of three runs per cell; standard errors smaller
than the markers); $f(q)$ is shown as a comparator, the truncated sets
do not satisfy the assumptions of the main-text Theorem~2.
(b)~Bootstrap distributions of $\hat\gamma$: primary family-stratified
cluster bootstrap (outline, dotted 95\% limits) and independent
resampling of runs (filled). (c)~Bootstrap distribution of the
family-norm reduction under the primary scheme; acceptance region
shaded.}
\label{fig:test}
\end{figure}

\paragraph*{Interpretation.} Once family fixed effects are included,
$\gamma$ is identified by within-family variation of $\ctil$, almost all
of it in the ring family, where $\ctil$ rises between $q=0.55$ and $0.65$
and is then constant at~1 while the advantage of the ring family over the
others shrinks. The negative coefficient records this co-variation; it is
evidence against the mediation hypothesis in this design, not evidence
that one-step containment lowers $M$. In Region~I the ring truncated
sets (contiguous arcs) give $M$ within $2\times10^{-4}$ of $f(\beta/n)$ at
$q\ge0.70$; the other families lie below $f$, which we do not explain.

% ============================================================
\section{Conjectures, open problems and limitations}
\label{sec:S8}
% ============================================================

\begin{conjecture}
\label{conj:upper}
In Region~II of the ring, $\E[M]\le f(q)$ for all $n$ and admissible $r$.
\end{conjecture}
\emph{Evidence.} Exact rationals for the six odd $n\le25$
(Table~\ref{tab:smexact}) and floating-point solutions of all 17
Region-II cases with $n\le33$ (main-text Table~II), with deviations from
$-0.0009$ to $-0.0195$, increasing monotonically with $q$. No coupling or
monotonicity proof is known.

\begin{openq}[Unequal balls]
Prove or refute $q^\ast\le\bar q$ under R1$\wedge$R2
(Remark~\ref{rem:het-numerics}).
\end{openq}

\begin{openq}[Fluctuations]
The variance of $M$ is not fixed by the affine closure, because it
involves pairwise overlaps $|B_j\cap B_k|$ of balls.
\end{openq}

\begin{openq}[Family dependence]
The difference between the configuration-model and 2D families in
Protocol~\ref{prot:reg} is unexplained; candidate causes include residual
ball-size heterogeneity, expansion properties and fractional-radius
effects.
\end{openq}

\begin{heuristic}[Fragmentation]
Below R2, minorities may fragment; no low-dimensional closed Markov
description in terms of shape variables is known.
\end{heuristic}

\paragraph*{Limitations.} (1)~The closed form requires two states, R1,
R2 and equal balls; Proposition~\ref{prop:het} relaxes the last.
(2)~Regions~II and~III have no closed-form characterization.
(3)~Claims~\ref{cl:B4}--\ref{cl:B5} are not used. (4)~The annealed size
dependence (Table~\ref{tab:sizescan}) is computed, not proved, and
Heuristic~\ref{heur:map} is a heuristic. (5)~The regression test uses one
graph per (family, $n$) and truncated interaction sets.

% ============================================================
\section{Validation ledger}
\label{sec:S9}
% ============================================================

\begin{table}[h]
\caption{Numerical checks cited in the main text and their evidence
files (paths relative to the repository root~\cite{repo}).}
\label{tab:ledger}
\small
\begin{ruledtabular}
\begin{tabular}{p{5.4cm}p{3.6cm}p{6.0cm}}
Check & Outcome & Evidence\\
\colrule
Ring containment formulas vs enumeration & exact, 192 cases & \path{results/validation/containment_bruteforce.json}\\
Exact chain, Region II, six cases, two solvers & agree; rationals certified & \path{results/validation/exact_independent.json}; \path{validation_results/T5.json}\\
Exact chain, Region I & equals $f(q)$ & \path{validation_results/T1.json}\\
Simulator vs exact chain, Region II & $|\Delta|\le4.6\times10^{-5}$ & \path{results/validation/sim_independent.json}\\
Simulator vs Theorem~\ref{thm:band}, Region I (ring, 2D) & $|\Delta|\le2.1\times10^{-5}$ & \path{results/validation/sim_independent.json}\\
Cycle law, $\E[K]$, $\mathrm{CV}_T^2$ & $|z|\le1.9$ & \path{results/validation/sim_independent.json}\\
Newborn identity & within binomial error & \path{validation_results/T6.json}\\
Region II, $n\le33$; annealed version & all below $f$; annealed below quenched & \path{scripts/region2_results.txt}\\
Proposition~\ref{prop:het}, 60 ball systems & $|\Delta|\le1.2\times10^{-15}$ & \path{scripts/hetero_balls_check.py}\\
Annealed size scan, $n\le3200$; random map & Table~\ref{tab:sizescan} & \path{scripts/annealed_size_scan.py}\\
Regression test & executed; criteria not met & \path{results/tables/reanalysis/}\\
\end{tabular}
\end{ruledtabular}
\end{table}

\bibliography{refs_pre}

\end{document}
